\section{Related Work} 
\label{sec:background}

%We describe two lines of research related to our work of discovering novel domains and intents from unlabeled text.

%\smallskip

%\noindent \textbf{Domain and Intent Detection:}
%Prior work on user intent detection encompasses two avenues: asynchronous, written communication (forums, blogs, tweets) and synchronous dialog from dialog agents such as Microsoft LUIS\footnote{https://www.luis.ai/}, Google Dialogflow\footnote{https://dialogflow.com/}, and Amazon Alexa\footnote{https://developer.amazon.com/alexa}. Supervised and semi-supervised learning models based on linguistic and sentiment features have been used to model racial intent on Tumblr~\cite{agarwal2017characterizing}, and purchase intent in social posts~\cite{gupta2014identifying,wang2015mining} and in discussion forums~\cite{chen2013identifying}. %Wang et al~\cite{wang2015mining} proposed a graph-based optimization approach for semi-supervised classification of tweets with a purchasing intent. Recent approaches have drawn upon progress in CNNs, RNNs, and transformer based language models to improve intent detection~\cite{xu2013convolutional,mesnil2015using,kim2016intent,liu2016attention,kim2017onenet,cai2017cnn,goo2018slot,xia2018,shivakumar2019spoken,castellucci2019multi} and domain detection~\cite{kim2018efficient,kim2018joint,zhang2018joint}. 
%Cai et al~\cite{cai2017cnn} used hierarchical clustering to learn a taxonomy of intent classes, and applied a hybrid CNN-LSTM model to classify the intent of medical queries. Performing slot filling jointly with intent detection has improved the performance of both tasks~\cite{kim2016intent,liu2016attention,zhang2016joint,wang2018bi}.  
Intent detection has been successfully performed in the literature via machine learning based approaches~\cite{tur2011sentence,jeong2008triangular,kim2016intent,ravuri2015comparative}, as well as deep learning models~\cite{sun2016weakly,xia2018,shivakumar2019spoken,castellucci2019multi,lin2019deep}. %\cite{ravuri2015comparative} performed a comparative study of different neural network architectures considering only lexical information of the utterance as features for intent detection. 
In recent years, intent detection has been jointly done with slot filling to improve its performance~\cite{xu2013convolutional,bhargava2013,liu2016attention,zhang2016joint,kim2017onenet,goo2018slot,wang2018bi}. However, the above approaches either require sufficient labeled data for each domain and intent, or some prior knowledge about the new intents to be discovered. \mname{} seeks to eliminate these restrictions.


%\smallskip

%\noindent \textbf{Open Classification:}
%Prior work classifies input instances into seen or unseen (during training) classes~\cite{scheirer2014probability,bendale2015towards,fei2016breaking,shu2017doc,xu2019open}. 
%Domain and intent detection 
Methods were proposed to recognize input texts belonging to novel domains~\cite{kim2018joint} and novel intents~\cite{lin2019deep,vedula2019towards}. However, unlike \mname{}, they do not jointly detect both novel domains and novel intents, cannot detect the number of novel domains and intents, and do not discover a taxonomy of the potential domain and intent types within user utterances. %likely to contain a novel domain or intent. 
Deep clustering networks~\cite{hsu2015neural} and Bayesian non-parametric models~\cite{dundar2012bayesian,akova2012self} were proposed to identify mixture components or clusters in data. But unlike \mname{}, the clusters identified cannot automatically be mapped to unique unseen categories. 
Our work %being completely unsupervised 
is also different from semi-supervised clustering techniques~\cite{fu2016semi,akova2012self,yi2015efficient,yi2013semi,hsu2015neural}, which lever small amounts of labeled data and require the classes to be discovered to be known beforehand. Having said that, we show the gains achieved by a semi-supervised version of \mname{}, in case prior knowledge is available about the novel intents (Table~\ref{tab:supervision}).